% Vector workflow, identity unit and inferential-boundary diagram.
\begin{figure}[tb]
\centering
\resizebox{\linewidth}{!}{%
\begin{tikzpicture}[
  >=latex,
  node distance=0.45cm,
  every node/.style={font=\footnotesize},
  auditbox/.style={draw=black!70,rounded corners=1.5pt,fill=black!3,
      inner sep=5pt,minimum height=1.50cm,text width=2.70cm,align=center},
  link/.style={->,line width=.55pt,draw=black!75}
]
\node[step] (release) {Public release\\50 distinct IDs\\2,000 processed trials};
\node[auditbox,right=of release] (hash) {Frozen EEG pipeline\\exact array + label\\candidate fingerprint};
\node[auditbox,right=of hash] (group) {Byte-confirmed overlap\\ID 4 / ID 5 full alias\\49 content groups};
\node[auditbox,right=of group] (cross) {Same 49 held groups\\old / purged source fit\\old / purged selection};
\node[auditbox,right=of cross] (effect) {Outcome audit\\reporting weight vs\\paired fit / selector};
\draw[link] (release) -- (hash);
\draw[link] (hash) -- (group);
\draw[link] (group) -- (cross);
\draw[link] (cross) -- (effect);
\end{tikzpicture}%
}
\caption{Unit-of-analysis and outcome-attribution workflow.
A candidate hash match is confirmed using numerical content and event
labels; its connected component defines a conservative split unit.
The group map uses already accessible public EEG signals and therefore
does not constitute a blind, zero-target-information discovery.
The matched four-cell comparison changes \emph{source fitting}
and \emph{source-selected regularisation} while retaining
the same 49 target contents; only the last stage accesses held
target outcomes. No stage resolves clinical participant identity.}
\label{fig:p2-identity-workflow}
\end{figure}
